% Exact table extracted from the pinned manuscript; needs booktabs/amsmath.
\begin{table}[t]
\centering
\caption{CRPS of hourly-only and combined second-scale attenuation.}
\label{tab:slow_fast_crps}
\small
\setlength{\tabcolsep}{4pt}
\renewcommand{\arraystretch}{1.1}
\resizebox{\linewidth}{!}{%
\begin{tabular}{@{}lrrrrrr@{}}
\toprule
Dataset & Hours & Hourly only & Combined & Reduction (\%) &
$\Delta$ (dB) & 95\% CI for $\Delta$ \\
\midrule
Local validation & 343 & 0.1413 & 0.1139 & 19.38 & -0.0274 & [-0.0294, -0.0256] \\
Local test & 687 & 0.1469 & 0.1190 & 18.96 & -0.0279 & [-0.0298, -0.0261] \\
UK pooled & 5593 & 0.2891 & 0.2474 & 14.41 & -0.0417 & [-0.0437, -0.0398] \\
Chilbolton & 3382 & 0.3514 & 0.3074 & 12.51 & -0.0439 & [-0.0463, -0.0417] \\
Sparsholt & 2211 & 0.1938 & 0.1556 & 19.69 & -0.0381 & [-0.0402, -0.0362] \\
\bottomrule
\end{tabular}
}
\par\smallskip
\begin{minipage}{\linewidth}\footnotesize
CRPS is in dB and uses identical valid seconds for both methods,
with equal weight per hour. The hourly-only distribution is the
frozen TabM $k=10$ point estimate; the combined distribution adds
32 generated fluctuations. $\Delta=C_{\mathrm{combined}}-
C_{\mathrm{hourly}}$; negative values favor the combination.
Relative reduction is calculated from unrounded mean scores.
Confidence intervals are paired date-cluster percentile bootstrap
intervals (2000 draws, seed 20261024), with both UK sites clustered
jointly by UTC date. These are intervals for the mean score
difference, not predictive bands for individual seconds.
\end{minipage}
\end{table}
